Earlier this year, a team of astronomers reported their discovery of a galaxy
with much less dark matter than the galaxy evolution model predicted (van
Dokkum et al.\ 2018, Nature). This galaxy, named NGC 1052-DF2, is located
close to the elliptical galaxy NGC 1052 ($D = 20$\,Mpc from the Sun) in the
sky. The shape of NGC 1052-DF2 resembles an ellipse with semi major axis ($a$)
of $22.6''$ and $\frac{b}{a} = 0.85$. Half of the total light from the galaxy
comes from within this ellipse and the mean surface brightness within the
ellipse is about 24.7\,mag\,arcsec$^{-2}$.

\begin{description}
\item[(a)] Calculate the total apparent magnitude of this galaxy.
\item[(b)] The team suggested the galaxy is a companion of NGC 1052.
  Determine the total mass of stars in NGC 1052-DF2, assuming it has a mass to
  light ratio $\left(\frac{M/M_\odot}{L/L_\odot}\right)$ of 2.0.
\item[(c)] The team identified 10 globular clusters in NGC 1052-DF2 with a
  mean galactocentric distance of $78.4''$. They also measured the velocity
  dispersion of these clusters to be not more than 8.4\,km/s. Estimate the
  dynamical mass of this galaxy. For simplicity, assume the mass distribution
  in the galaxies is uniform and is spherically symmetric.
\item[(d)] This discovery was challenged by other groups (Kroupa et al.,
  Nature, 2018, Truijlo et al., MNRAS, 2018), who claimed that NGC 1052-DF2 is
  not a satellite of NGC 1052, and it is located at a much smaller distance to
  us. Show why a smaller distance would weaken the assertion of the dark
  matter deficiency in NGC 1052-DF2.
\end{description}
